% Results section (draft 14 Nov 2026). All numbers copied from docs/results_ledger.md; do not retype from memory.
% Figures: results/figures/fig_rmse_by_depth.pdf, results/figures/fig_argo_thermocline_by_year.pdf (copy into the Overleaf project).

\section{Results}
\label{sec:results}

\subsection{Agreement with GLORYS}
On the headline Test~A period the frozen ensemble reproduces GLORYS with a pooled RMSE of 0.671$^\circ$C on the deep set, against 0.969$^\circ$C for the seasonal climatology and 0.684$^\circ$C for the per-cell linear regression (Table~\ref{tab:main}); the corresponding skill is $+0.52$. On Test~B the RMSEs are 0.729, 1.068 and 0.746$^\circ$C. The ensemble improves on the linear regression by $1.8\%$ (95\% interval $[-2.9,-0.7]$) on Test~A and $2.4\%$ ($[-3.2,-1.6]$) on Test~B, and on XGBoost by $1.6\%$ and $1.7\%$ (intervals excluding zero). A single U-Net does not beat the linear regression (0.690 and 0.745$^\circ$C); averaging members improves it by $2.7\%$ and $2.2\%$. The demeaned ensemble is worse than every baseline against GLORYS (0.743 and 0.810$^\circ$C).

\begin{table}[t]
\centering
\caption{Pooled RMSE ($^\circ$C) over 15 depths on the deep set, against GLORYS12 and against Argo profiles. Argo is assimilated by GLORYS and is not independent of it. Best learned model in bold per column.}
\label{tab:main}
\begin{tabular}{@{}lcccc@{}}
\hline
 & \multicolumn{2}{c}{vs GLORYS} & \multicolumn{2}{c}{vs Argo} \\
Model & Test A & Test B & Test A & Test B \\
\hline
Climatology & 0.969 & 1.068 & 1.009 & 1.106 \\
Linear (SLA+SST) & 0.684 & 0.746 & 0.861 & 0.932 \\
XGBoost & 0.682 & 0.741 & 0.855 & 0.920 \\
Single U-Net & 0.690 & 0.745 & 0.848 & 0.917 \\
U-Net ensemble & \textbf{0.671} & \textbf{0.729} & 0.834 & 0.909 \\
Ensemble, demeaned SLA & 0.743 & 0.810 & \textbf{0.754} & \textbf{0.848} \\
\hline
GLORYS12 & -- & -- & 0.677 & 0.783 \\
\hline
\end{tabular}
\end{table}

\begin{figure*}[t]
\centering
\includegraphics[width=\textwidth]{figures/fig_rmse_by_depth.pdf}
\caption{RMSE against GLORYS12 by depth on the deep set for Test~A (a) and Test~B (b), and (c) the RMSE change of the U-Net ensemble relative to the linear regression with 95\% block-bootstrap intervals (negative is better). The depth axis is a square-root scale.}
\label{fig:depth}
\end{figure*}

The improvement is concentrated in the thermocline (Fig.~\ref{fig:depth}). Relative to the linear regression it reaches $-5.6\%$ at 75~m on Test~A and about $-4\%$ at 50 and 150~m on Test~B, with intervals excluding zero at 50--150~m in both periods. At 0--10~m the ensemble is \emph{worse} than the linear regression in both periods (by $2.7$--$4.8\%$ of the RMSE, at most $0.02^\circ$C). Below 300~m all models are at the level of the climatology (RMSE $0.33$--$0.39^\circ$C); the ensemble's gains there are statistically detectable but at most $0.01^\circ$C, and we do not interpret them as skill.

By region (Arabian Sea, Bay of Bengal, Andaman Sea; boxes fixed before the test periods were scored), the gain over the linear regression is clear in both periods only in the Bay of Bengal ($-5.6\%$ and $-3.8\%$). In the Arabian Sea, which holds $65\%$ of the deep-set cells, it is not distinguishable from zero on Test~A ($-0.5\%$, $[-1.8,+0.7]$) and $-1.9\%$ on Test~B. The advantage seen in the Andaman Sea on the validation window ($-12.5\%$) did not reproduce ($+0.6\%$ on Test~A, $-3.4\%$ on Test~B). By season on Test~B, the gain is $-3.7\%$ in DJF, $-3.5\%$ in MAM and $-1.8\%$ in SON, and absent in JJA ($-0.9\%$, $[-2.0,+0.1]$); Test~A is too short (16 months) for seasonal intervals.

\subsection{Agreement with Argo profiles}
Compared profile by profile with Argo (Table~\ref{tab:main}), GLORYS itself has an RMSE of 0.677 and 0.783$^\circ$C on Test~A and Test~B, and the models are further away, as expected for estimates trained to reproduce GLORYS. The frozen ensemble agrees with Argo to 0.834 and 0.909$^\circ$C against 1.009 and 1.106$^\circ$C for the climatology, and is closer than the linear regression by $3.1\%$ ($[-5.3,-1.3]$) and $2.5\%$ ($[-3.8,-1.3]$), but the advantage is confined to 75--200~m ($-5.5\%$ in both periods): at 5--50~m it is worse than the linear regression ($+3.8\%$ and $+2.7\%$).

The ranking between the two ensembles reverses with the reference. Against Argo the demeaned ensemble is the closest learned model (0.754 and 0.848$^\circ$C), $9.7\%$ ($[-12.8,-7.7]$) and $6.7\%$ ($[-8.5,-4.9]$) closer than the frozen ensemble, whereas against GLORYS it is the worst. Fig.~\ref{fig:argo} shows why the reference matters: in the thermocline (75--200~m) GLORYS is warmer than Argo in every year (by $0.22$ to $0.56^\circ$C), and the models trained on it inherit this bias and amplify it in some years (the frozen ensemble reaches $+0.79^\circ$C in 2024). The demeaned ensemble has smaller yearly biases (root-mean-square $0.30$ against $0.48^\circ$C for the frozen ensemble) but their sign changes between years ($-0.30$ to $+0.51^\circ$C). After removing each year's mean bias, the two ensembles have the same RMSE in the thermocline ($1.024$ and $1.026^\circ$C), so the demeaned ensemble's advantage is a reduction of bias, not better pattern skill; both are $3.6$--$3.7\%$ better than the linear regression in that centred measure. Argo coverage is thin in the Andaman Sea (8 profiles in Test~B), where skill cannot be checked against profiles.

\begin{figure*}[t]
\centering
\includegraphics[width=\textwidth]{figures/fig_argo_thermocline_by_year.pdf}
\caption{Mean bias (a) and RMSE (b) against Argo at 75--200~m by calendar year (validation and test windows; training years excluded; 2020 and 2021 span two windows). Tick labels give the number of Argo profile-depth values.}
\label{fig:argo}
\end{figure*}

\subsection{Which inputs the networks use}
Integrated Gradients on 19 validation days shows that the networks rely on SST near the surface ($63\%$ of the absolute attribution at 0~m, $50\%$ at 30~m) and on sea level below about 50~m ($72\%$ at 50~m, $83\%$ at 150~m); the same-day fields carry about half of the attribution and the 7-, 14- and 30-day lags about $19\%$, $13\%$ and $14\%$. Attribution describes reliance, not correctness: sea level also carries most of the attribution at 500--1000~m, where no model has skill.
